% ================================================================
%  第 2 章 巡游一部作品——或孩子与母亲 (后半, 原书 p.54-71, 边码 P38-P65)
%  源文: _work/src1/task-ch02b.md
%  说明: 原书注 49、注 60 的注文在源文中缺失, 正文保留裸上标
% ================================================================

% ---------------- p.54 ----------------
\origpage{P38} 还在此前不久,我们关于这些上同调不变量的构想,已因让·勒雷的工作而大大丰富、深刻更新——那工作是战争期间他在德国当俘虏时继续做的,时在四十年代前半。其决定性的创新想法,是空间上的(阿贝尔)层,勒雷给每个层配上一列相应的“上同调群”(所谓“以该层为系数的”上同调群)。这就好像那根一直用来“丈量”空间的、可靠的老式“上同调标尺”,忽然被增殖为多得难以想象的众多新“标尺”,各种尺寸、形状、质地都有,每一根都贴合于所论的那个空间,每一根都就它向我们交出精确无比的信息,而且是唯有它才能给出的信息。这正是我们对各类空间的研究方式发生深刻转变的主导想法,也肯定是本世纪出现过的最关键的想法之一。多亏让–皮埃尔·塞尔后来的工作,勒雷的想法很快结出最初的果实——在其出现后的头一个十年里,拓扑空间(尤其是与上同调紧密相关的所谓“同伦”不变量)理论有了一次令人瞩目的重新启动;而所谓“抽象”代数几何也有另一次同样根本的重新启动(靠塞尔那篇奠基性文章 FAC,1955年发表)。我自1955年起在几何方面的工作,正与塞尔这些工作相接续,因而也与勒雷那些创新想法相接续。

\section{(13) topos——或双人床}

勒雷引入的层的视角与语言,让我们以一种新的光照去看各类“空间”与“簇”。然而它们并未触及空间概念本身,只是让我们以新的眼睛更精细地把握这些人人早已熟悉的传统“空间”。可后来表明:这个空间概念,并不足以承载那些最本质的“拓扑不变量”——正是它们表达着“抽象”代数簇(例如韦伊猜想所适用的那些)的“形状”,乃至一般“概形”(旧日簇的推广)的形状。对于那场期待中的“数与量的联姻”来说,这就像一张分明过于狭窄的床:充其量只有未来夫妇中的一个(即新娘)能勉强挤进去安身,却从来容不下两个!而那尚待发现的“新原则”——好让仁慈的仙女们许下的婚约得以圆房——不是别的,也正是未来夫妇所缺的那张宽敞的“床”;而在此之前,竟没有人哪怕察觉到这一点。

这张“双人床”,随着 topos(广义空间)这个想法出现了(像被魔杖一点……)。这个想法在同一种拓扑直觉里,既容纳了体现着连续量世界的传统(拓扑)空间,也容纳了那些不知悔改的抽象代数几何学家所谓的“空间”(或“簇”),还容纳了无数其他类型的结构——它们此前似乎无可挽回地被铆在“不连续”或“离散”的聚集体的“算术世界”上。

正是层的视角做了那沉默而可靠的向导,那把有效(且绝非秘密)的钥匙,
% ---------------- p.55 ----------------
\origpage{P39} 毫不迟疑、毫不绕弯地把我领进洞房,领到那张宽大的婚床前。这床确实宽得(像一条宽广、宁静、深不见底的河……),宽到——
\begin{quote}
“国王的全部马匹\\
都能一齐来饮……”
\end{quote}
——正如一支老歌里唱的;那歌你肯定也唱过,至少听人唱过。而头一个唱这歌的人,对 topos 那种隐秘的美、那种宁静的力量,感受得比我那些博学的学生、比我旧日的任何朋友都要真切……

钥匙是同一把,无论在最初那套临时性的进路里(经由“site”这个十分方便、却并非内蕴的概念),还是在 topos 的进路里。现在我想试着描述一下 topos 这个想法。

试想一个给定(拓扑)空间上所有层所构成的“集合”,或者,愿意的话,由所有这些用来“丈量”它的“标尺”所构成的那套了不起的装备\fnote{44}{}{(给数学家看的)严格说来,这里说的是集合的层,而不是勒雷为了构造“上同调群”而引入、作为最一般系数的阿贝尔层。顺便说,我相信自己是头一个系统地使用集合的层的人(从1955年起,在堪萨斯大学写的那篇《具有结构层的纤维空间的一般理论》里)。}。我们把这一“集合”或这套“装备”看作带有它最显然的那种结构——如果可以说的话,那结构“一眼就看得见”;也就是所谓“范畴”结构。(非数学的读者不必因为不知道这个术语的技术含义而发慌。下文他完全用不着它。)正是这种可称为“(所论空间上的)层范畴”的“丈量上层建筑”,从此将被看作“体现”着空间最本质的东西。这在(“数学上的健全常识”看来)完全正当,因为事实是:人们可以凭一个拓扑空间所关联的这个“层范畴”(或这套“丈量装备”),把它整个地“重新构造”出来\fnote{45}{}{(给数学家看的)严格说来,这只对所谓“索伯(sobre)空间”成立。不过这类空间几乎包括了我们平常遇到的全部空间,尤其包括分析学家们所钟爱的一切“分离”空间。}。要核实这一点只是个简单的练习——当然,得先把问题提出来……有了这一条,就足以确信:只要我们出于某种理由觉得合适,从此便可以“忘掉”最初那个空间,只保留、只使用与之关联的那个“范畴”(或那套“装备”);它将被看作所要表达的“拓扑(或空间)结构”最恰当的化身。

正如数学中常有的事,我们在这里(靠“层”或“上同调标尺”这个关键想法)成功地把某个概念(这里是“空间”)用另一个概念(这里是“范畴”)表达了出来。每一次,当人们发现可以把一个概念(表达某一类情境)译成另一个概念(对应另一类情境)时,我们对这两个概念的理解都会同时加深——因为分别附着于这一个或那一个的特定直觉,会出人意料地汇流到一起。于是,一个“拓扑”性质的情境(由某个给定空间体现),在这里被译成了一个“代数”性质的情境(由某个“范畴”体现);或者,愿意的话也可以说:由空间体现的“连续”,被这个“代数”性质的范畴结构“译”出、“表达”出来了——甚至直到那“不连续”或“离散”性质的外包装。

但这里还有更多的东西。这些概念中的头一个,即空间概念,在我们看来多少是个“极大”的概念——它已经如此一般,以至很难想象还能找到一个仍然“合理”的扩张。然而,事实是:在镜子的另一面\fnote{46}{}{这里说的“镜子”,像《爱丽丝镜中奇遇》里那面一样,是把一个空间放在它前面、便给出该空间之“像”的那面镜子;像就是与之关联的“范畴”,被看作空间的一种“替身”,“在镜子的另一面”……},我们从拓扑空间出发所遇到的这些“范畴”(或“装备”),性质非常特别。它们确实享有一组特征极强的性质\fnote{47}{}{(给数学家看的)这里说的主要是我在范畴论中以“正合性”为名引入的那些性质(同时还引入了现代范畴论中一般的归纳极限与投射极限概念)。见《论同调代数中的若干要点》,《东北数学杂志》,1957年(第119–221页)。},使它们看上去像是其中最简单那一个的各式“仿制品”——那最简单的,就是从缩成单独一点的空间出发所得到的那个。话虽如此,一个推广了传统拓扑空间的“新型空间(或 topos)”,将被径直描述为一个“范畴”:它未必来自某个寻常空间,却仍然具备这样一个“层范畴”所应有的全部好性质(当然,这些性质是一次性地明文规定好的)。

% ---------------- p.56 ----------------
\origpage{P41} 这就是那个新想法。它的出现,可以看作下面这个其实近乎孩子气的观察的结果:在一个拓扑空间里,真正要紧的绝不是它的“点”或点的子集\fnote{48}{}{例如,人们可以造出非常“大”的 topos,它只有一个“点”,甚至根本没有“点”!},也不是点与点之间的邻近关系之类,而是这个空间上的层,以及它们所构成的范畴。总而言之,我不过是把勒雷最初的想法推到了它最终的结论——推完之后,再跨出那一步。

如同层这个想法本身(源自勒雷),如同概形这个想法,如同一切要来撞翻某种积习已深的看法的“大想法”,topos 这个想法也自有一种令人不知所措之处:它那么自然,那么“显然”,那么单纯(可以说,到了天真乃至简陋的极限,甚至到了“犯傻”的地步)——就是那种特别的性质,常让我们脱口而出:“哦,原来就这么点东西!”语气里一半失望,一半嫉妒;也许还外加一层“古怪”“不正经”的弦外之音,凡是因为出人意料的过度简单而叫人摸不着头脑的东西,往往就落得这评语。这也许是在提醒我们,想起童年那些久已埋没、久已否认的日子……

\section{(14) 空间概念的变异——或气息与信仰}

概形概念构成了“代数簇”概念的一次大扩张,因此它把我前人所传的代数几何从里到外更新了一遍。topos 概念则构成了空间概念的一次意想不到的外延,更确切地说,一次蜕变。由此,它承载着拓扑学、并在此之外还有几何学获得类似更新的许诺。而且,从现在起它已经在新几何的兴起中起了关键作用(主要通过由它而生的 $\ell$-进上同调与晶体上同调这两个主题,并通过它们,在韦伊猜想的证明中起作用)。像它的姐姐(而且几乎是孪生姐姐)一样,它具备任何一次丰产的推广所必备的两个互补特征,如下。

% ---------------- p.57 ----------------
\origpage{P42} 第一,新概念不太宽:也就是说,在新的“空间”里(人们更愿意叫它们“topos”,免得冒犯一些敏感的耳朵\fnote{49}{}{“topos”这个名字是(连同“topologie”或“topologique”之名一起)选定的,意在暗示:它正是拓扑直觉所适用的“那个对象”。凭这个名字唤起的丰富心象之云,应当把它看作大致相当于“(拓扑)空间”一词,只是更强调这个概念“拓扑”上的特殊性。(于是,有“向量空间”,但在另行通知之前,并没有“向量 topos”!)看来应当把两个说法一并保留,各存其特殊性。}),那些对旧日好空间来说司空见惯的、最本质的“几何”直觉与构造\fnote{50}{}{这些“构造”里,尤其包括一切熟悉的“拓扑不变量”的构造,上同调不变量也在其中。至于后者,我在已经提到的那篇文章(“东北”1955)里把该做的都做了,好让它们对任何“topos”都有意义。},大体上都能搬过来。换句话说,对这些新对象,我们手头有整套丰富的意象与联想、概念以及至少某些技术,而这些东西以前只限于老式对象。

第二,新概念同时又足够宽,能把许多情境包进来——这些情境在此之前并不被认为会产生“拓扑–几何”性质的直觉,而那正是过去只留给寻常拓扑空间的直觉(而且自有其道理……)。

从韦伊猜想的眼光看,这里的关键在于:新概念确实足够宽,使我们能给每个“概形”配上这样一个“广义空间”或“topos”(称为所论概形的“平展 topos”)。这个 topos 的某些“上同调不变量”(全是些“傻乎乎”的东西!)当时看来颇有希望提供“所需要的东西”,使这些猜想获得全部意义,而且(谁知道呢!)也许还能提供证明它们的办法。

正是在我此刻正在写下的这些纸页里,我有生以来第一次(作为数学家)有闲暇去述说(哪怕只是对自己述说)我数学作品中全部的主旋律与重大指导性想法。这使我更好地领会这些主题、以及它们所体现的种种“视角”,在统一它们、并生出它们的那幅宏大几何图景中各自的位置与分量。正是通过这项工作,新几何最初那强劲兴起中两个要害的创新想法才在充分的光照下显现出来:概形的想法,与 topos 的想法。

这两个想法里,如今在我看来更深的是第二个,即 topos 的想法。假如五十年代末我没有挽起袖子,日复一日固执地干下去,一干十二个长年,做出一件精细与力量都臻完美的“概形工具”——那么我几乎无法想象,在此后的一二十年里,别人最终会一直忍住不去引入那个明明非引入不可的概念(哪怕是被迫的),也不会好歹搭起几间“预制件”的破棚屋,聊胜于我这番有心一石一石拼合、亲手竖起的那一座座宽敞舒适的宅邸。然而反过来,在近三十年的数学舞台上,我看不出还有谁会有这份天真,或这份单纯,去(替我)跨出那另一步——那一步比什么都关键,它引入了 topos 这个如此孩子气的想法(哪怕只是“site”那个想法)。就算这个概念已经有人殷勤地奉上,连同它似乎暗藏的那点怯生生的许诺——我也看不出还有谁,无论旧友还是学生,会有那份气息,尤其是那份信仰,把这卑微的想法\fnote{51}{}{(给数学家读者)我说“把这卑微的想法做到底”,指的是把平展上同调当作通往韦伊猜想的一条进路这个想法。正是在这句话的启发下,我于1958年发现了 site 概念;而这个概念(或与之极相近的 topos 概念)以及平展上同调的形式体系,则是在我的推动下于1962至1966年间发展起来的(有几位合作者协助,他们将在适当的地方谈到)。\par{}我说“气息”和“信仰”,指的是“非技术”性质的品质,而在这里,它们在我看来正是本质的品质。在另一个层面上,我还可以加上我所谓的“上同调的嗅觉”,也就是当初在我身上发展起来的、用于建造上同调理论的那种嗅觉。我原以为把它传给了我的上同调学生们。离开数学界十七年之后回头一看,我发现他们当中没有一个人把它保留下来。}一直做到底(它表面上如此微不足道,而目标看来无限遥远……):从最初结结巴巴的起步,直到“平展上同调之掌握”完全成熟——在随后那些年里,它最终在我手中化成了这个。

% ---------------- p.58 ----------------
\section{(15) 国王的全部马匹}

\origpage{P43}\origpage{P44} 是的,河很深;我童年的水又宽又静,在那片我以为早已离开的国土上。国王的全部马匹都能在那里从容地、喝个痛快地一齐来饮,也饮不尽它!那水来自冰川,炽热如那些遥远的积雪,又有着平原黏土的柔和。我刚说到其中一匹马,一个孩子把它牵来饮水,它久久地喝了个够。我还看见另一匹来饮了一会儿,说不定是循着同一个孩子的足迹来的——可那一回没耽搁多久。一定是有人把它赶走了。也就这样了,可以说,再没有别的。然而我看见平原上流浪着无数干渴的马群——就在今天早上,它们的嘶鸣还把我从床上拽起来,时辰还早得不像话,而我已是快六十的人,喜欢清静。没办法,我只好起床。看着它们瘦成一副副饿马的骨架,我心里难受,而好水明明不缺,青草地也不缺。可是看来,有人对这片我记忆中殷勤好客的国土施了恶意的魔法,封住了通往那些丰沛之水的路。又或许这是本地的马贩子设的局,好把价钱压下来,谁知道呢?又或许这是个再也没有孩子牵着马去饮水的国度,马儿们只好渴着,因为再没有一个孩子找得到通往那条河的路……

\section{(16) Motif——或心中的心}

topos 这个主题出自概形主题,就在概形出现的那一年——但在广度上它远远超出了母题。正是 topos 主题(而非概形主题)才是那张“床”,或那条“深河”:几何与代数、拓扑学与算术、数理逻辑与范畴论、连续的世界与“不连续”或“离散”结构的世界,都在那里结亲。如果说概形主题好比新几何的心脏,那么 topos 主题就是它的外壳,或它的居所。它是我所构想的最宽广的东西,为的是用同一种充满几何回响的语言,精细地把握那些彼此相距最远的情境所共有的某种“本质”——那些情境来自数学事物这个辽阔宇宙的这一片或那一片。

然而 topos 这个主题远没有像概形主题那样走运。我在《收获与播种》里多处谈到此事,这里不是细说这个概念所遭遇的种种离奇变故的地方。不过新几何的两个主旋律主题都出自 topos 主题,那是两个互补的“上同调理论”,它们各自都是为着给韦伊猜想提供一条进路而构想出来的:平展(或“$\ell$-进”)主题,与晶体主题。前者在我手中化成了 $\ell$-进上同调工具,如今它已显出是本世纪最强有力的数学工具之一。至于晶体主题,在我离开之后被压缩到一种近乎隐秘的存在,最终(在需要的压力下)于1981年6月被发掘出来,在聚光灯下,用一个化名,而情形比围绕 topos 的那些还要离奇。

正如所预料的,$\ell$-进上同调工具成了建立韦伊猜想的基本工具。我自己证明了其中一大摞,而最后一步是在我离开三年后由皮埃尔·德利涅出色地完成的,他是我那些“上同调”学生中最出色的一位。

另外,约在1968年,我曾提炼出韦伊猜想的一个更强的版本,而且尤其是一个更“几何”的版本。这些猜想一直“带着”(如果可以说的话!)一层看来不可约化的“算术”色彩,而这些猜想的本意恰恰是要通过“几何的”(或“连续的”)中介,来表达和把握“算术”(或“离散的”)\fnote{52}{}{(给数学家看的)韦伊猜想从属于一些“算术”性质的假设,尤其因为所论的簇必须定义在有限域上。从上同调形式主义的观点看,这导致要给与这种情形相关联的 Frobenius 自同态一个特殊位置。在我的进路里,那些关键性质(“广义指标定理”一类)涉及任意的代数对应,并且对事先给定的基域不作任何算术性质的假设。}。
% ---------------- p.59 ----------------
\origpage{P45}\origpage{P46} 在这个意义上,我所提炼的那个版本,比韦伊本人的版本更“忠于”“韦伊的哲学”——那种从未写出、也很少被说出的哲学;在过去四十年几何的非凡兴起中,它也许是主要的、未明言的动机\fnote{53}{}{不过,在我1970年离开之后,出现了一股十分明显的反挫,具体表现为一种相对停滞的局面;我在《收获与播种》的字里行间不止一次有机会谈到它。}。我的改写,要点在于提炼出某种“精华”:在被称作“抽象”代数簇的框架内,那仍然应当有效的东西——出自对“寻常”代数簇有效的经典“霍奇理论”\fnote{54}{}{这里“寻常”的意思是:“定义在复数域上”。霍奇理论(所谓“调和积分”理论)是复数代数簇背景下已知上同调理论中最强的一个。}。我把这些著名猜想的这个全新版本(就代数闭链而言)称为“标准猜想”。

在我心里,这是在发展出 $\ell$-进上同调工具之后,朝这些猜想迈出的新的一步。但同时,而且尤其是:它也是通往那个在我看来仍是我在数学中引入的最深主题的可能进路之一\fnote{55}{}{这是最深的主题,至少在我数学家活动“公开”的时期(1950至1969年)里是这样,也就是说直到我离开数学舞台为止。我认为自1977年起发展起来的远阿贝尔代数几何与伽罗瓦–泰希米勒理论这一主题,深度与之相当。}:Motif 的主题(它本身也生于“$\ell$-进上同调主题”)。这个主题好比概形主题的心脏或灵魂,是它最隐蔽、最躲开目光的部分;而概形主题本身又处于新图景的核心。而标准猜想\fnote{56}{}{(给代数几何读者)必要时,这些猜想需要重新表述。更详细的评论见《工地巡览》(ReS IV 第178号注,第1215–1216页)以及《信念与知识》第769页的页底注(ReS III 第162号注)。}里提炼出的少数关键现象,可以看作构成了 Motif 主题的终极精华,看作这个在一切主题中最精微者的“气息”,看作新几何这颗“心中的心”的“气息”。

大致是这么回事。我们已经看到,对给定的素数 $p$,(尤其为着韦伊猜想)知道如何为“特征 $p$ 的(代数)簇”构造“上同调理论”,这有多重要。而那条著名的“$\ell$-进上同调工具”恰恰提供了这样一种理论,甚至提供了无穷多个不同的上同调理论——对每个异于特征 $p$ 的素数 $\ell$,各有一个。这里显然又有一件“缺失的理论”,它对应的正是 $\ell$ 等于 $p$ 的情形。为补这个缺,我专门设想了另一种上同调理论(前面已经提到过),叫作“晶体上同调”。另外,在 $p$ 为无穷这个重要的情形里,还另有三种上同调理论\fnote{57}{}{(给数学家读者)这些理论分别对应于 Betti 上同调(用超越方法定义,借助把基域嵌入复数域)、霍奇上同调(塞尔定义)和 de Rham 上同调(我定义);后两者已可上溯到五十年代(而 Betti 上同调则上溯到上个世纪)。}——而且谁也不能保证人们迟早不会再去引入新的上同调理论,其形式性质全都类似。于是,与寻常拓扑学里的情形相反,这里摆在我们面前的是一大堆多得令人惶惑的不同上同调理论。人们有一种很清楚的印象:在某种起初还相当模糊的意义上,所有这些理论都该“归于同一”,它们“给出同样的结果”\fnote{58}{}{(给数学家读者)例如,若 $f$ 是代数簇 $X$ 的自同态,它诱导出上同调空间 $H^{i}(X)$ 的一个自同态,那么后者的“特征多项式”应当整系数,并且不依赖于所选取的那个特定上同调理论(例如 $\ell$-进理论,$\ell$ 可变)。对于一般的代数对应也一样,只要 $X$ 假定是真且光滑的。可悲的真相(它也让人看出,在我离开之后,特征 $p>0$ 的代数簇的上同调理论处于多么可悲的荒废状态)是:这件事至今仍未得到证明,甚至在 $X$ 是射影光滑曲面、$i=2$ 这个特殊情形里也没有。事实上,据我所知,在我离开之后还没有人肯费心关注这个关键问题——它正是那种看来从属于标准猜想的典型问题。时尚的法令是:唯一值得注意的自同态是 Frobenius 自同态(德利涅用现成的家伙什把它单独处理掉了……)。}。

% ---------------- p.60 ----------------
\origpage{P47} 正是为了表达不同上同调理论之间这种“亲缘”直觉,我提炼出了与一个代数簇相关联的“Motif”概念。用这个词,我想暗示:在借助先天地可能的全部上同调理论、与簇相关联的这一大堆不同上同调不变量之下,有一个“共同的动机”(或“共同的理由”)。这些不同的上同调理论,就好比同一个“基本动机”(称为“Motif 上同调理论”)的若干不同的主题展开,各自带着自己的“速度”、“调性”和“调式”(“大调”或“小调”);而这个基本动机,同时又是所有这些不同的主题化身(也就是说,所有这些可能的上同调理论)中最根本、最“精细”的一个。这样,与一个代数簇相关联的 Motif,便构成了“终极的”、“头等的”上同调不变量;其他一切(与各种可能的上同调理论相关联者)都从它推出,如同种种音乐化身,或种种“不同的实现”。簇的“上同调”的一切本质性质,都早已能在相应的 Motif 上“读出”(或“听出”),于是,那些特殊化的上同调不变量(例如 $\ell$-进或晶体不变量)上我们所熟悉的性质与结构,便只是 Motif 内部性质与结构的忠实映照\fnote{59}{}{(给数学家读者)另一种看域 $k$ 上 Motif 范畴的方式,是把它想象成 $k$ 上有限型分离概形范畴的一种“阿贝尔包络范畴”。于是,与这样一个概形 $X$ 相关联的 Motif(或“$X$ 的 Motif 上同调”,我记作 $H_{mot}^{*}(X)$)便显出是 $X$ 的一种阿贝尔化的“化身”。这里的关键在于:正如一个代数簇 $X$ 可以有“连续变化”(它的同构类因而依赖于连续的“参数”,或“模”),与 $X$ 相关联的 Motif、或更一般地一个“变化的”Motif,也同样可以有连续变化。这是 Motif 上同调的一个侧面,它与一切经典上同调不变量的情形形成鲜明对照,包括 $\ell$-进不变量——唯一的例外是复数代数簇的霍奇上同调。\par{}这就让人想见,“Motif 上同调”是一个精细得多的不变量,比传统上纯粹拓扑的不变量更紧地圈定了 $X$ 的“算术形状”(如果我敢冒险用这个说法)。在我的 Motif 图景里,这些 Motif 构成一条极其隐蔽、极其精致的“绳索”,把一个代数簇的代数–几何性质,与由它的 Motif 所体现的“算术”性质联结起来。后者可以被看作一个就其精神本身而言“几何”性质的对象,但在其中,那些从属于几何的“算术”性质可以说被“裸露”出来了。\par{}这样,在我眼中,Motif 便是人们至今能够与一个代数簇相关联的最深刻的“形状不变量”,除了它的“Motif 基本群”之外。这两个不变量在我看来,都像是一种尚待描述的“Motif 同伦型”的“影子”(我在《工地巡览——或工具与眼光》那条注里顺便说了几句(ReS IV 第178号,见工地5(Motif),尤其是第1214页))。我觉得,正是后面这个对象,才该是任意代数簇那种难以捉摸的“算术”(或“Motif”)“形状”直觉最完美的化身。}。

以上用音乐比喻的非技术语言说出的,是一个仍然具有孩子般单纯的想法之精华,它既精细又大胆。我是在那些我认为更紧迫的基础工作之余发展这个想法的,名之曰“Motif 理论”或“Motif 的哲学(或‘瑜伽’)”,时为1963–1969年。这是一个结构丰富得令人着迷的理论,其中很大一部分至今仍是猜想性的\fnote{60}{}{这些年里,但凡有人愿意听,我就把自己的 Motif 图景讲给他听,却始终不肯费事把这件事白纸黑字地发表出来(服务于众人的其他任务,我并不缺)。这就使后来我的某些学生能更从容地劫掠,而在我那些对情形一清二楚的旧友们的温情注视之下。(参见下页页底的注。)}。

% ---------------- p.61 ----------------
\origpage{P48}\origpage{P49} 我在《收获与播种》里多次谈到这套“Motif 的瑜伽”,它格外让我牵挂。这里不是回头重述我在别处所说的地方。只需说:那些“标准猜想”再自然不过地从这套 Motif 瑜伽中流出。同时,它们也为 Motif 概念的一种可能的成形构造提供了一个进路。

这些猜想,在我看来——今天在我看来仍然如此——是代数几何中最基本的两个问题之一。这个问题,以及另一个同样关键的问题(所谓“奇点消解”),至今都还没有解决。然而,第二个问题在今天仍像一百年前一样,显得是一个显赫而令人生畏的问题;而我曾有幸提炼出的那一个,却被时尚的专断法令(就在我离开数学舞台之后的那些年里,而且正如 Motif 主题本身一样\fnote{61}{}{事实上,这个主题在1982年(晶体主题之后一年)被发掘出来,这一回用的是它原来的名字(而且形式狭窄,只在基域特征为零这一种情形里),却没有说出工匠的名字。这只是一个例子,还有若干个:某个概念或主题在我离开之后随即被当作格罗滕迪克式的幻影埋掉,随后十年或十五年里又被我的某些学生一个个发掘出来,带着谦逊的骄傲,而且(还用得着再说明吗)不提工匠……})归为格罗滕迪克式的可爱胡闹。不过我又在预先说后面的事了……

\section{(17) 发现母亲——或两个坡面}

说实话,我为了建立韦伊猜想而对它们本身所作的思考,一直是零星的。那片已开始在我面前展开、我努力去细看、去把握的全景,在广度和深度上都远远超出一次证明的假定所需,甚至远远超出这些著名猜想最初所能让人瞥见的一切。随着概形主题与 topos 主题的出现,一个意想不到的新世界骤然打开了。“那些猜想”在其中确实占据中心位置,有点像一片幅员辽阔、行省无数的帝国或大陆的首都;而那些行省大多数与这个辉煌显赫的地方只有最遥远的关联。我从未需要对自己明说,但我知道自己从此是一项大任务的仆役:去探索这个巨大而未知的世界,把握它的轮廓直到最远的边界;也去四处走遍,以顽强而有条理的心思清点那些最近、最易到达的行省,为它们绘制忠实而精确得一丝不苟的地图,连最小的村落、最小的茅屋也要有它的位置……

主要是后面这项工作吸走了我大部分精力——那是一项耐心而庞大的基础工作,只有我一人看得清楚,尤其是只有我一人“用五脏六腑感觉得到”。从1958年(概形主题与 topos 主题接连出现的那一年)到1970年(我离开数学舞台的那一年),占去我时间远为最多的正是它。

而且我常常咬着嚼子忍耐:被这样拖住,像被一块黏人的重物坠着,没完没了地做那些(一旦看清了要害之后)在我看来更近乎“后勤”而非向未知冲刺的活儿。我得不断地按住那股向前冲的冲动——那是拓荒者或探险者的冲动,他出发去发现、去探索一个个未知而无名的世界,那些世界不停地呼唤我,要我去认识它们、为它们命名。那股冲动,以及我投在它上面的精力(几乎像是偷来的!),始终少得可怜。

然而我心底很清楚:正是这份从“任务”那里(可以说)偷来的精力,才是最稀有、最精微的本质——我作为数学家的“创造”,首先正落在那里:落在那份专注的凝神上,去把握那尚未诞生之物
% ---------------- p.62 ----------------
\origpage{P50}\origpage{P51} 在一个温暖而不竭的孕育母体中幽暗、无形、潮湿的褶皱里初现的形状与轮廓的痕迹;它似乎在呼唤我,要成形、要化身、要诞生……在发现的工作里,这份专注的凝神、这份炽热的关切是一种本质的力量,正如太阳的温暖之于埋在沃土里的种子的幽暗孕育,之于它们卑微而奇迹般地迎光绽放。

在我作为数学家的工作中,我看得最清楚的是这两股力量或冲动,它们同样深,而(我觉得)性质不同。为了说它们,我用了“建造者”的形象,以及“拓荒者”或“探险者”的形象。把两者并排放着,我忽然觉得它们都太“阳”,太“男性”,甚至“大男子气”!它们带着神话般高傲的回响,或“重大时刻”的回响。它们肯定是被我身上旧日那种“英雄式”的创造工作观(那超级阳刚的观点)的残迹所激发的。就这样,它们给出一种浓重着色的、甚至僵硬的、“立正”式的看法,去看一个远比这流动、远比这谦卑、远比这“简单”的现实——一个活生生的现实。

然而在这股“建造者”的男性冲动里(它似乎不断把我推向新工地),我同时也清楚辨认出恋家者的冲动:深深依恋“那个”家的人的冲动。首先,那是“他的”家,是“亲近者”的家——是一个亲密的活体之所在,他觉得自己是它的一部分。只是到了后来,随着被感受为“亲近”的圈子逐渐扩大,它才也成为一座“众人的家”。而在这股“造房子”的冲动里(就像人们“做”爱那样……),也有,而且首先有,一份柔情。还有与那些材料接触的冲动:一件一件地塑造它们,带着爱怜的用心,而只有通过这种充满爱的接触才能真正认识它们。而当墙立起来、梁和屋顶架好之后,还有一种深深的满足:一间一间地安置,看着在这些厅堂、房间和小角落里,渐渐建立起这座活着的房子的和谐秩序——美、好客、宜于居住。因为房子,首先并且隐秘地在每个人心里,也是母亲——是环绕我们、庇护我们的东西,既是庇护所又是安慰;而且也许(更深一层,甚至就在我们正把它整个造出来的时候)我们自身也正是从那里来的,是它在我们的出生之前那些永远被遗忘的时光里庇护了我们、哺育了我们。那也是“怀抱”。

而刚才为了越过“拓荒者”这个显赫名号、为了把握它所遮盖的那个更隐蔽的现实,自然而然浮现出来的形象,也同样剥去了一切“英雄”腔调。那里出现的,仍然是母性的原型意象——哺育的“母体”及其无形而幽暗的劳作……

这两股在我看来“性质不同”的冲动,最终比我以为的更接近。两者都属于一种“接触的冲动”,把我们带向“母亲”的相遇:带向那位既体现着切近的、“已知”之物,又体现着“未知”之物的她。委身于这一股或那一股冲动,都是“重新找到母亲”。都是重新接上与切近者、与“或多或少已知者”的接触,同时也接上与“遥远者”、与那“未知”却同时被预感着、即将被认识者的接触。

这里的差别是色调上的、配比上的,不是性质上的。当我“造房子”时,占上风的是“已知”;当我“探索”时,是未知。这两种发现的“方式”,或者更确切地说,同一个过程或同一项工作的这两个方面,不可分割地连在一起。两者都本质,而且互补。在我的数学工作里,我辨认出在这两种进路之间、更确切地说在其中一个占上风的时刻(或时期)与另一个占上风的时刻(或时期)之间,有一种不断的往返运动\fnote{62}{}{我在这里就数学工作所说的话,对“静观”的工作也同样成立(《收获与播种》里到处都会谈到它)。我几乎不怀疑:这是出现在一切发现工作中的东西,包括艺术家(譬如说,作家或诗人)的工作。我在这里描述的两个“坡面”,也可以看作:一个是表达及其“技术”要求的坡面,另一个是接受(种种知觉与印象)的坡面,后者靠一份专注的凝神而变成灵感。两者在工作的任何时刻都在场,而且在其中一个占上风的“时节”与另一个占上风的“时节”之间,总有这种不断的“往返”运动。}。但同样清楚的是:在任何时刻,两种方式都在场。

% ---------------- p.63 ----------------
\origpage{P52} 当我建造、布置,或者清理、打扫、归置时,定调的是工作中“阳”的或“男性”的“方式”或“坡面”。当我摸索着探索那不可把握者、那无形者、那无名者时,我便是自己存在中“阴”的或“女性”的那一面。

对我而言,问题绝不是要贬低或否弃我天性中的这一面或那一面,两者都本质——那建造并生育的“男性”,和那孕育、并庇护缓慢幽暗的怀胎的“女性”。我“是”两者——“阳”与“阴”,“男人”与“女人”。但我也知道,创造过程里最精细、最精微的本质,在“阴”的、“女性”的那一坡面一边——那谦卑、幽暗、往往其貌不扬的坡面。

工作中始终(我相信)对我施以最强魅惑的,正是这一坡面。可当时的共识却鼓励我把最明澈的那部分精力投到另一坡面,投到那在可触的“产品”中化身并自我确证的坡面——且不说那些产品还是完成、收尾了的,轮廓分明,以凿过的石头般的明证来担保自己的实在……

拉开距离看,我很清楚这些共识是怎样压在我身上的,也很清楚我是怎样“承受这重量”的——柔顺地承受!我工作中“构思”或“探索”的那一部分,直到我离开的时候都还被压缩在最低限度,好吧。然而,在这样回望我作为数学家的作品时,有一点以惊人的明显浮现出来:使这部作品有其本质与力量的,恰恰是这个如今被轻视、甚至被嘲笑或居高临下地鄙夷的坡面:那是“想法”的坡面,甚至是“梦”的坡面,绝不是“结果”的坡面。在这些纸页里,我试着界定我对我那个时代的数学所贡献的最本质的东西,用的是一种拥抱森林、而不在几棵树上流连的眼光——我看到的不是一串“大定理”的荣誉榜,而是一把活生生的、丰产想法的扇子\fnote{63}{}{并非我的作品里缺少人们所谓的“大定理”,其中还包括解决了别人提出、而我之前无人能解的问题的定理。(我在“上涨的海……”那条注(ReS III 第122号)第554页的页底注(***)里回顾了其中几个。)但正如我在这趟“巡游”开头(“视角与眼光”那一站,第6节)已经强调过的,这些定理对我来说,只有在一个大主题的滋养性语境里才获得全部意义,而那主题是由某个“丰产想法”开启的。于是它们的证明就像从源头、毫不费力地从承载它们的那个主题的本性、“深度”中流出——正如河水之浪仿佛从它水域的深处温柔地生出,没有断裂,也不费力。我在已经引过的那条注“上涨的海……”里,用别的意象表达了完全类似的意思。},它们全都汇聚到同一个广阔的图景里。

\section{(18) 孩子与母亲}

当这篇“前言”开始变成一趟穿越我数学作品的巡游,还加上我那段关于“(血统纯正的)继承人”和“(无可救药的)建造者”的小小品时,这个没写成的前言也开始有了一个名字:它该叫“孩子与建造者”。接下来几天里越来越清楚:“孩子”和“建造者”是同一个人物。于是这名字变得更简单了:“孩子—建造者”。天哪,这名字倒不缺少气派,而且完全对我的胃口!

可只见反思表明:这位高傲的“建造者”,或者(更谦逊地说)那个玩造房子游戏的孩子,只是那个著名的“玩耍的孩子”的一副面孔,而他有两副。还有另一个孩子:他爱探索东西,爱到处乱掏,爱钻进沙堆,钻进泥泞无名的淤滩,
% ---------------- p.64 ----------------
\origpage{P53}\origpage{P54} 钻进最不可能、最荒唐的角落……大概是为了掩人耳目(哪怕只是瞒过我自己……),我起初给他起了个火亮的名字“拓荒者”,后来又换成“探险者”——这名字更接地气,却仍罩着一圈威望的光晕。简直要问:在“建造者”和“拓荒者–探险者”之间,哪一个更男性、更诱人!正面还是反面?

可再凑近一看,我们这位无畏的“拓荒者”原来是个女孩(我却乐意把她打扮成男孩)——她是水洼、雨、细雨和夜的姐妹,沉默,几乎看不见,因为她总退到阴影里——那个总被遗忘的人(人们不对她做出一副嘲弄的样子就已经不错了……)。而我也真有办法,一天又一天地把她忘掉——可以说,是双倍地忘掉:我起初只愿意看见那个男孩(那个玩造房子的孩子……)——即便后来我终于忍不住,到底还是看见了另一个,我看到的她,也仍然是个男孩……

至于我那趟巡游的好名字,这下就完全站不住了。那是个全“阳”的、全“大男子”的名字,一个瘸腿的名字。要让它站得不歪,就得把另一个也写进去。可怪的是,“另一个”并没有真正的名字。勉强贴得上的只有“探险者”,但那还是个男孩的名字,没办法。这里的语言是个婊子,它甚至不让你察觉就把你套住,显然跟祖传的偏见串通一气。

也许可以用“造房子的孩子与探索的孩子”敷衍过去。把这一点留在不言之中:一个是“男孩”,另一个是“女孩”,而这是同一个男孩–女孩的孩子,他在建造时探索,在探索时建造……可是昨天,除了这观照与探索、命名与建造的阴阳两坡面之外,事情的又一面还显现了出来。

宇宙、世界,乃至万有,归根到底是陌生而遥远的东西。它们并不真正与我们相干。在我们最深处,求知的冲动并不是把我们带向它们。吸引我们的,是它们可触而切近的化身,最近、最“肉感”的那个,充满深沉的回响与奥秘——那化身与我们血肉之躯的起源相混同,也与我们族类的起源相混同——也是那自古以来就等待着我们、沉默着准备接纳我们的她,“在路的另一端”。冲动从她那里涌出,欲望的道路从她那里伸展——从她,母亲,从那位生我们如同生这世界的她;而那些道路把我们带向与她的相遇,朝她奔去,为的是一次次折返,消融在她之中。

% ---------------- p.65 ----------------
\origpage{P55} 就这样,在一趟不期而至的“巡游”的路上拐个弯,我意外地重又遇到一个曾让我觉得亲切、又被我有些遗忘了的寓言——孩子与母亲的寓言。可以把它看作一个关于“生命寻求自身”的寓言。或者在更谦卑的个体存在层面上,看作一个关于“存在者寻求事物”的寓言。

这是一个寓言,也是一种祖辈经验的表现,那经验深植于心灵之中——是滋养深层创造层的原初象征中最有力的一个。我相信在其中认出了以原型意象的远古语言表达出来的东西:人身上创造力的气息本身,它鼓动着人的血肉与精神,无论在它最卑微、最短暂的表现里,还是在最辉煌、最持久的表现里。

这“气息”,正如使它化身的血肉形象一样,是世上最卑微的东西。它也是最脆弱的、最被人忽视、最遭人轻贱的东西……

而这气息在你一生中的种种遭际,不是别的,正是你的冒险,你一生中“认识”的冒险。表达它的那个无言的寓言,就是孩子与母亲的寓言。

你是孩子,出自母亲,被庇护在她之中,以她的力量为食。孩子从母亲那里奔出——那至近的、熟识的母亲——去与母亲相遇:那无边的、永远未知而满是奥秘的母亲……

\begin{center}
“巡游一部作品”终
\end{center}

\subsection*{尾声:看不见的圈子}
\tocsubsub{尾声:看不见的圈子}

\section{(19) 死亡是我的摇篮(或三个娃娃给一个垂死者)}

直到五十年代末 topos 视角出现为止,空间概念的演变在我看来都是一种本质上“连续”的演变。它似乎平顺地、没有跳跃地延续着,从欧几里得对我们周围空间的理论化开始,从希腊人传下的几何学开始——那几何学致力于研究生活在这个空间里的某些“图形”(直线、平面、圆、三角形等)。诚然,数学家或“自然哲学家”构想“空间”\fnote{64}{}{我写尾声时最初的打算,是附一份极简的梗概,略述其中某些“深刻变化”,并显出我在其中看到的这种“本质连续性”。我放弃了,免得把这趟巡游拉得太长——它已经比预定的长多了!我想在《反思》第四卷拟议的“历史评注”里再谈,那一回是为数学家读者写的(这完全改变了阐述的任务)。}的方式发生过深刻的变化。但这些变化在我看来全都属于一种本质的“连续性”——它们从未把(像所有人一样)依恋着熟悉意象的数学家,突然抛进一片陌生的地方。它们就像一个人身上那些也许深刻、却是渐进的变化:我们从这人的童年就认识他,从他会走路一直看到他成年、完全成熟。有些漫长平淡的时期里,变化觉察不到;另一些时期里,也许又很剧烈。但即便在生长或成熟最强烈的时期,即便我们有几个月、甚至几年没再见到他,任何时刻都不可能有丝毫怀疑、丝毫犹豫:我们重又见到的,仍然是他,一个熟识的、亲近的存在,哪怕容貌已经变了。

另外我想可以说:到本世纪中叶,这个亲近的存在已经大大地老了——像一个终于筋疲力尽、磨损殆尽的人,被一大堆他毫无准备的新任务压垮。也许他甚至已经安然死去,却没人费心去记下一笔、验明此事。“大家”还都做出一副在活人家里忙活的样子,弄得几乎就像他确实还活得好好的。

那么,请想想这对家里的常客是多么难堪:那位可敬的老人原本僵直地端坐在扶手椅里,忽然却看见一个身量不过三个苹果高的壮实小崽子在那儿撒欢,还顺口宣称——毫不带笑,仿佛理所当然——空间先生(从现在起,您尽可随意把那“先生”二字去掉……)就是他!要是他至少有点像家里的相貌也好啊,说不定是个私生子,谁知道……可一点也没有!一眼看去,没有任何地方让人想起那位大家那么熟悉(或者以为熟悉……)的老空间父亲,那位大家无论如何都确信(这总该是最起码的……)是永恒的父亲……

% ---------------- p.66 ----------------
\origpage{P56}\origpage{P57} 这就是那著名的“空间概念的变异”。这就是我至少从六十年代初起就不得不“看见”、把它当作明摆着的事的东西,却一直没机会把它对自己明确说出来,直到我写下这些行字的此刻。而忽然,只凭这个形象的说法、只凭它立刻唤起的联想之云,我以新的清晰看见:传统的“空间”概念,以及与之紧密相关的“簇”的概念(各类簇,尤其是“代数簇”),在我来到这一带的时候,已经老得那么厉害,简直就像已经死了……\fnote{65}{}{这个断言(有人会觉得它武断)得“加点盐”来听。它的有效程度,不多不少,正如下面那个断言(我稍后也把它算在自己账上):力学的“牛顿模型”(地上的或天上的)在本世纪初已是“奄奄一息”,那时爱因斯坦赶来救驾。事实是,即便在今天,物理学里大多数“寻常”情形下,牛顿模型都完全够用,而去找相对论模型才是发疯(鉴于测量所允许的误差幅度)。同样,数学里许多情形下,旧的熟悉的“空间”与“簇”概念仍然完全够用,不必去找幂零元、topos 或“温和结构”。但在两种情形里,对于尖端研究中越来越多出现的语境,旧的概念框架都已变得无法表达哪怕最“寻常”的情境。}。我还可以说:正是随着概形视角(连同它的子嗣\fnote{66}{}{(给数学家看的)在这个“子嗣”里,我尤其算上形式概形、各类“重数”(特别是概形重数或形式重数),最后还有所谓“刚性解析”空间(由泰特引入,依循我提供的一个“总施工”,灵感来自 topos 这个新概念,同时也来自形式概形的概念)。这份清单当然远不完备……},外加一万页基础工作为其代价)、随后是 topos 视角的接连出现,一个不声不响的危机局面终于被解开了。

在刚才那幅图景里,其实不该说是一个小崽子,作为一次突然变异的产品,而该说是两个。而且这两个孩子之间有一种无可否认的“家族相似”,尽管他们一点也不像那位故去的老人。再者,凑近看还可以说:概形这小家伙,好比是已故的空间父亲(别名“各类簇”)与 topos 这小家伙之间的一个“亲缘环节”\fnote{67}{}{另外,对这两个娃娃,还该再加上第三个更小的,它出现在不那么宽厚的年代:就是“温和空间”这小崽子。正如我在别处指出过的,它没拿到出生证明,而我还是在完全不合法的状态下把它算进了我有幸引入数学的十二个“主旋律主题”之中。}。

\section{(20) 瞥一眼对门的邻居}

这局面,我觉得与本世纪初爱因斯坦的相对论理论出现时的局面非常接近。那时有一个更触目的概念死胡同,表现为一个突然的、看来无法解决的矛盾。理所当然,那个要来重整混沌的新想法,是一个孩子般单纯的想法。值得注意的是(而且符合一个再重复不过的剧本……):在所有那些忽然紧张起来、拼命想“保住家什”的聪明人、显要人、大人物当中,没有一个人想到这个想法。非得是一个无名的年轻人,刚从(说不定)学生的阶梯教室板凳上下来,跑来(也许还为自己那份大胆有点不好意思……)向他那些显赫的前辈解释该怎么做才能“拯救现象”:只要不再把空间与时间分开就行了\fnote{68}{}{这样描述爱因斯坦的想法,当然有点简略。在技术层面上,需要弄清的是在新时空上放什么结构(不过这一点其实已经“飘在空中”了,有麦克斯韦的理论和洛伦兹的想法)。这里本质性的一步不是技术性质的,而是“哲学”性质的:意识到对于相距遥远的事件而言,同时性概念没有任何实验上的实在性。就是这个“孩子气的观察”,这句“可皇帝没穿衣服!”,让人跨过了那道著名的“限定一个宇宙的、专横而看不见的圈子”……}!从技术上说,当时一切都已具备,足以让这个想法绽开并被接受。而爱因斯坦那些前辈的体面也正在于此:他们确实懂得接纳新想法,没有过多训斥。这是一个迹象,说明那还是个伟大的时代……

% ---------------- p.67 ----------------
\origpage{P58} 从数学角度看,爱因斯坦的新想法平淡无奇。但从我们对物理空间的构想来看,那却是一次深刻的变异,一次突然的“背井离乡”。这是自2400年前欧几里得提炼出物理空间的数学模型以来的第一次此类变异;那模型自古至今被所有物理学家和天文学家(包括牛顿)原封不动地拿来供力学之用,用以描述地上与天上的力学现象。

爱因斯坦这个最初的想法后来大大深化了,化成一个更精微、更丰富、更柔韧的数学模型,借助了已有数学概念这座丰富的武库\fnote{69}{}{这里主要是指“黎曼簇”的概念,以及这种簇上的张量计算。}。随着“广义相对论”,这个想法扩展为一幅物理世界的宏大图景,以同一道目光拥抱亚原子的极小微观世界、太阳系、银河与遥远的星系,以及电磁波在时空中行进——那时空每一点都被其中的物质所弯曲\fnote{70}{}{这个模型区别于空间与时间的欧几里得(或牛顿)模型、也区别于爱因斯坦最初那个模型(“狭义相对论”)的最触目的特征之一,是时空的整体拓扑形状始终不确定,而不是由模型本身的性质强制规定。这种整体形状究竟是什么,在我看来(作为数学家)是宇宙学中最迷人的问题之一。}。在宇宙学与物理学的历史上,这是第二次、也是最后一次(继三个世纪前牛顿那第一次大综合之后),出现一幅用数学模型的语言表达的、统一的宏大图景,涵盖宇宙中全部物理现象。

不过,爱因斯坦这幅物理宇宙的图景,后来也被事态淹没了。需要解释的“全部物理现象”,从本世纪初起就有时间不断扩充!出现了大量物理理论,各自以或多或少的成功,去解释那一大堆“被观察的事实”这座巨大杂物堆里有限的一捆。而人们至今还在等那个大胆的孩子,他会在玩耍中找到那把新钥匙(如果真有那么一把的话……),那个梦想中的“蛋糕模型”,肯于一下子“走通”,把所有现象一齐拯救……\fnote{71}{}{人们把这样一种假想的理论称为“统一理论”,它终将把前面谈到的众多局部理论“统一”起来、调和起来。我觉得,那项尚待着手的基本思考,将要在两个不同的层面上展开。\par 1°)一个“哲学”性质的思考,关于“数学模型”这个概念本身——为一个现实片段而设的数学模型。自牛顿理论成功以来,物理学家的一条默许公理便是:存在一个数学模型(甚至是唯一一个模型,或“那个”模型),能完美地表达物理实在,没有“脱节”,没有含糊。这个统治了两个多世纪的共识,像是毕达哥拉斯“万物皆数”那活生生的洞见留下的一块化石残迹。也许这就是新的“看不见的圈子”,取代了从前的形而上学圈子,用来限定物理学家的宇宙(而“自然哲学家”这一族类看来已彻底灭绝,被计算机那一族高奏凯歌地取代了……)。只要肯停下来哪怕一瞬,就很清楚:这个共识的有效性绝非显而易见。甚至有非常严肃的哲学理由,导致先天地怀疑它,或至少为它的有效性划定极严格的界限。现在正是时候,把这个公理交给一番严格的批判,甚至“证明”——超出一切可能的怀疑——它没有根据:不存在唯一严格的数学模型,能说明至今编目的全部所谓“物理”现象。\par 一旦“数学模型”这个概念本身、以及这样一个模型的“有效性”概念(在测量所允许的某些“误差幅度”之内)被令人满意地界定清楚,“统一理论”的问题、或至少“最优模型”(意义尚待明确)的问题,才终于会被清楚地提出来。同时,对于选择这样一个模型所带有的(也许是必然带有的)任意性程度,人们大概也会有更清楚的想法。\par 2°)在我看来,只有在这样的思考之后,提炼一个比前人的模型更令人满意的明确模型这一“技术”问题,才获得全部意义。那时也许正是时候,摆脱物理学家的第二条默许公理;这条公理可上溯到古代,深植于我们感知空间的方式本身:那就是空间与时间(或时空)的连续本性,也就是“物理现象”在其中展开的那个“场所”的连续本性。\par 想必已是十五、二十年前,我翻着构成黎曼全集的那本薄薄的书,被他“顺带”的一句评论打动了。他在那里指出:空间的终极结构完全可能是“离散的”,而我们关于它所形成的“连续”表象,也许只是对一个更复杂的实在的一种(久而久之也许过度的)简化;对人类精神来说,“连续”比“不连续”更容易把握,因此它便充当了我们把握不连续的一种“近似”。这话出自一位数学家之口,洞察力令人吃惊,而且当时物理空间的欧几里得模型还从未受到质疑;从严格的逻辑意义上说,传统上反倒是“不连续”被用作通向连续的技术进路。\par 何况近几十年来数学的发展表明,连续结构与不连续结构之间的共生,远比本世纪前半人们所想象的更为紧密。无论如何,要找到一个“令人满意的”模型(或者必要时,一组这样的模型,尽可能令人满意地“彼此衔接”),无论它是“连续”的、“离散”的还是“混合”性质的——这样的工作肯定要动用极大的概念想象力,以及一种炉火纯青的嗅觉,去把握并揭示新型的数学结构。这种想象力或“嗅觉”,我觉得是稀罕之物,不仅在物理学家中如此(那里爱因斯坦和薛定谔似乎是难得的例外),甚至在数学家中也如此(而这话我是完全有把握说的)。\par 总而言之,我预计,那期待中的更新(如果它还会来的话……)更可能出自一位骨子里的数学家、同时又熟知物理学的重大问题,而不是出自一位物理学家。但尤其是:那需要一个人具备“哲学上的开放”,才能抓住问题的症结。这个症结绝不是技术性质的,而是一个“自然哲学”的根本问题。})。

% ---------------- p.68 ----------------
\origpage{P59}\origpage{P60} 把我对我那个时代数学的贡献,与爱因斯坦对物理学的贡献相比较,这一念头是出于两个理由逼到我面前的:这两份作品都得益于我们对“空间”的构想发生的一次变异(一者取数学意义,另一者取物理意义);而且两者都采取一种统一性图景的形式,拥抱极大一批此前看来彼此分离的现象与情境。我在此看到他的作品\fnote{72}{}{我绝不声称自己熟悉爱因斯坦的作品。事实上,他的著作我一本也没读过,他的想法我只是通过道听途说、而且极为粗略地知道。不过我觉得自己还是看得出“森林”,尽管我从未费劲去细看他的任何一棵树……}与我的作品之间一种明显的精神亲缘。

这种亲缘,我觉得丝毫不被一种明显的“实质”差异所否定。正如我前面已经暗示的,爱因斯坦式的变异涉及物理空间概念,而爱因斯坦是从已知数学概念的武库里取用,从不需要扩充它,更不必颠覆它。他的贡献在于:在当时的已知数学结构中,挑出那些最适合\fnote{73}{}{关于“奄奄一息”这个形容词的评论,见前面的一条页底注(第55页的注)。}充当物理现象世界之“模型”的,用来取代前人传下的那个垂死的模型。在这个意义上,他的作品确实是一位物理学家的作品,并且在此之外,还是一位“自然哲学家”的作品——取牛顿及其同代人所理解的意义。这个“哲学”维度在我的数学作品里是没有的:我从未被引去追问那些在数学事物的宇宙中进行的“理想”概念构造,与物理宇宙中发生的现象(甚至心灵中经历的事件)之间可能有什么关系。我的作品是一位数学家的作品,他有意避开了“应用”(于其他科学)的问题,也避开了我工作的“动机”与心理根源的问题。而且是一位被其十分独特的天赋所驱使、不断扩充其技艺根基处概念武库的数学家的作品。就这样,我甚至没有察觉、仿佛在玩耍之间,被引去颠覆了对几何学家来说最根本的那个概念:空间概念(以及“簇”概念),也就是我们对几何对象所居住的那个“场所”本身的构想。

新的空间概念(像一种“广义空间”,但那些本该构成“空间”的点已经或多或少消失了)在实质上一点也不像爱因斯坦引入物理学的那概念(后者对数学家来说毫不令人迷失)。相反,比较倒该与薛定谔发现的量子力学去做\fnote{74}{}{据我理解(通过从各方传回的回声),人们一般认为本世纪物理学有过三次“革命”或大动荡:爱因斯坦的理论、居里夫妇对放射性的发现、以及薛定谔引入量子力学。}。在这种新力学里,传统的“质点”消失了,代之以一种“概率云”,它在此处与彼处的周围空间中疏密不同,取决于质点
% ---------------- p.69 ----------------
\origpage{P61}\origpage{P62} 位于该区域的概率。在这种新眼光里,人们清楚感到我们构想力学现象的方式发生了比爱因斯坦模型所体现的更深的一次“变异”——这次变异并不是简单地把一个在肩宽上略嫌狭窄的数学模型,换成一个类似但裁得更宽、更合身的模型。这一回,新模型与那些好的旧传统模型像得如此之少,以至连力学的大专家数学家也必定突然感到背井离乡,甚至迷失(或者愤慨……)。从牛顿力学走到爱因斯坦力学,对数学家来说,想必有点像从地道的普罗旺斯土话走到最时新的巴黎黑话。而走到量子力学,我想,那是从法语走到汉语。

而这些“概率云”——取代了昔日那些令人安心的物质粒子——奇异地让我想起充满 topos 的那些飘忽不定的“开邻域”:它们像一个个消散的幽灵,环绕着一些想象出来的“点”;而不肯驯服的想象力仍旧不顾一切地攀附着那些点……

\section{(21) “唯一者”——或孤独的馈赠}

这趟到“对门邻居”物理学家家里的小小远足,对一个(像大多数人那样)对数学家的世界一无所知、却肯定听人说起过爱因斯坦和他那著名的“第四维”、甚至听说过量子力学的读者,也许可以充当一个参照点。说到底,即便发明者当初并未料到他们的发现会具体化为广岛,后来又化为军备与(所谓)“和平”用途的原子竞赛,事实是:物理学上的发现有形地、几乎立刻地影响着一般人的世界。数学发现的影响,尤其在所谓“纯”数学(也就是说,不以“应用”为动机)里,则不那么直接,肯定也更难界定。譬如我就不知道,我对数学的贡献曾“服务”于什么,比如说造出个什么机器。事情如此,我当然无功可居,但这不妨碍它让我安心。一旦有了应用,就可以肯定:头一个扑上去的是军方(其后是警察)——至于工业(哪怕是所谓“和平”工业),也不见得总好多少……

当然,无论为我自己导航,还是为一位数学家读者导航,更该做的其实是在数学史本身之中找“参照点”来定位我的作品,而不是到别处去找类比。这些天我想过这事,以我对那段历史相当模糊的了解为限\fnote{75}{}{我还是个小鬼的时候,就对历史不太上心(地理也一样)。(《收获与播种》第五部分(只写了一部分)里,我曾“顺带”探到我觉得这种对历史的局部“阻塞”的深层原因——我相信近些年来这阻塞正在消解。)何况我的前辈们在“布尔巴基圈子”里所接受的数学教育,也没让事情好转——那里偶尔出现的历史引用少而又少。}。在“巡游”途中,我已经有机会提到一条数学家的“谱系”,他们的性情是我所认同的:埃瓦里斯特·伽罗瓦(Evariste Galois)、黎曼、希尔伯特。如果我对本门技艺的历史知道得更多,说不定我能把这条谱系往过去延伸得更远,或者在其中插入几个我只有耳闻的名字。让我吃惊的是:我不记得曾听说过(哪怕只是通过比我更熟悉历史的朋友或同行的暗示)除我之外还有哪个数学家,带来过一大串创新想法,而这些想法并不是彼此多少脱节的,而是像一个宏大统一图景的各个部分(物理与宇宙学中的牛顿和爱因斯坦、生物学中的达尔文和巴斯德就是这样的情形)。在数学史上,我只知道两个“时刻”诞生过一幅宏大格局的新图景。其一是2500年前古希腊数学作为我们今天所理解的科学诞生之时。另一个首先是十七世纪微积分诞生之时,那个时代以牛顿、莱布尼茨、笛卡尔等人的名字为标志。据我所知,在这两个时刻之一诞生的图景,都不是一个人的作品,而是一个时代的集体作品。

% ---------------- p.70 ----------------
\origpage{P63} 当然,从毕达哥拉斯和欧几里得的时代到十七世纪初,数学有时间改变面貌;同样,从十七世纪数学家创造的“无穷小计算”到本世纪中叶,也是如此。但据我所知,在这两个时期(一个两千多年,一个三百年)里发生的深刻变化,从未具体化、凝聚为一幅表达在某一部给定作品里的新图景\fnote{76}{}{写完这些行字几个小时后,我忽然想到:我在这里竟没有想到布尔巴基先生那部(集体)论著所力图呈现的当代数学的大综合。(在《收获与播种》第一部分里还会大量谈到布尔巴基小组。)我觉得这出于两个理由。\par 一方面,这一综合仅限于把一大批已经知道的想法与结果来一番“整理”,并不从中贡献它自己独创的新想法。若说其中有新想法,那大概是“结构”概念的一个精确数学定义;它确实成了贯穿整部论著的一条宝贵线索。但这个想法在我看来,更像一位聪明而富有想象力的词典编纂者的想法,而不像一种语言获得更新的要素——那种更新会带来对实在(这里是数学事物的实在)的新把握。\par 另一方面,从五十年代起,结构这一想法就被事态超越了:随着“范畴”方法骤然涌入数学中某些最有活力的部分,例如拓扑学或代数几何。(这样,“topos”概念就拒绝进入布尔巴基那只结构的“口袋”,那口袋确实窄得在肩宽上不够用!)布尔巴基固然是在完全知情的情况下决定不去趟这“浑水”,但这样一来,它也就放弃了最初的雄心:为整个当代数学提供基础与基本语言。\par 不过,它倒是固定下了一套语言,同时也固定下某种写作风格与处理数学的方式。这种风格起初是某种精神的(很不完整的)反映,那精神是希尔伯特活泼而直接的遗产。五十、六十年代,这种风格终于占了上风——好的方面如此,(更主要地)坏的方面也如此。近二十年来,它终于变成一种僵硬的“教规”,其所标榜的“严格”纯是门面,而当年鼓舞它的那种精神,看来已一去不返。},不像物理学和宇宙学里随大综合而发生的情形:牛顿的,然后是爱因斯坦的,在它们历史中的两个关键时刻。

看来,作为我心中诞生的那幅统一性宏大图景的仆役,我在数学史上从开端到如今都算“独一无二”。抱歉,我这样子像是要把自己抬举得超过了允许的分寸!不过,让我自己松一口气的是,我似乎辨认出一个潜在的(而且是天赐的!)兄弟。

刚才我已经有机会提到他,作为我那些“性情上的兄弟”中的第一位:埃瓦里斯特·伽罗瓦。在他短暂而闪电般的一生中\fnote{77}{}{埃瓦里斯特·伽罗瓦(1811–1832)死于一场决斗,年二十一岁。我记得他有好几部传记。我年轻时读过物理学家英费尔德写的一部传记体小说,当时给我很深的印象。},我相信辨认出了一幅宏大图景的开端——正是那“数与量的联姻”的图景,在一幅新的几何图景之中。我在《收获与播种》别处提到\fnote{78}{}{见《伽罗瓦的遗产》(ReS I,第7节)。},两年前我心中如何突然出现这个直觉:在当时最强烈地魅惑我的那项数学工作里,我正在“接续伽罗瓦的遗产”。这个直觉此后很少被提起,却有时间在沉默中成熟。我三周来对作品的回顾性思考,肯定又对此有所助益。如今我相信自己与一位过去数学家认出的最直接的谱系关系,正是把我与埃瓦里斯特·伽罗瓦连起来的那一条。不管对错,我觉得我一生中十五年发展起来、并在我离开数学舞台之后十六年里仍在我心中继续成熟、继续丰富的那幅图景——这图景也是伽罗瓦不可能不去发展的那一幅\fnote{79}{}{另外我深信,一个伽罗瓦会走得比我远得多。一方面因为他那份完全卓异的天赋(就我而言,并未分得这份天赋)。另一方面因为他大概不会像我这样,让大部分精力被那些没完没了的活儿分去,去把已经多少到手的东西逐一细致地整理成文……},
% ---------------- p.71 ----------------
\origpage{P64}\origpage{P65} 假如他处在我这个位置待在这一带,而没有早逝粗暴地截断一段壮丽的冲势。

肯定还有一个理由,促成我这种“本质亲缘”的感觉——这种亲缘不能归结为单单“数学性情”,也不能归结为一部作品的显著方面。在他的生活与我的生活之间,我还感到一种命运上的亲缘。诚然,伽罗瓦死得愚蠢,才二十一岁;而我呢,快六十了,还决心活得长长久久。可这并不妨碍:埃瓦里斯特·伽罗瓦生前,正如一个半世纪之后的我,在官方数学界里都是一个“边缘人”。就伽罗瓦的情形,浅看之下会觉得这种边缘性是“偶然的”,他不过是还没来得及靠自己的创新想法和作品“站住脚”。就我的情形,在我数学家的最初三年里,我的边缘性是由于我对数学界本身的存在一无所知(也许是故意的……)——那个我本该去面对的圈子;而自我离开数学舞台以来的十六年里,它是深思熟虑的选择的结果。肯定正是这个选择,招来了报复性的“无懈可击的集体意志”:要把我的名字的一切痕迹从数学中抹去,连同我所侍奉的那幅图景一起抹去。

但越过这些偶然的差异,我相信在这种“边缘性”里辨认出一个共同的原因,我觉得它是本质的。这个原因,我不把它看作历史境况,也不看作“性情”或“性格”上的特点(这些在我与他之间,大概就像人与人之间可能有的那么不同),当然更不在“天赋”的层面(伽罗瓦显然天赋惊人,而我比较起来只是平平)。如果真有一种“本质亲缘”,我看见它在一个远远更谦卑、远远更基本的层面上。

在我一生中,我曾在少数难得的场合感到这样的亲缘。也正因它,我觉得自己与另一位数学家“相近”,他是我的前辈:克劳德·谢瓦莱\fnote{80}{}{我在《收获与播种》里多处谈到克劳德·谢瓦莱,尤其是“与克劳德·谢瓦莱相遇——或自由与善意”那一节(ReS I 第11节),以及“向克劳德·谢瓦莱告别”那条注(ReS III 第100号注)。}。我要说的联系,是某种“天真”或“单纯”的联系,我前面有机会谈到过它。它表现为一种倾向(常不为周围人所赏识):用自己的眼睛去看事物,而不是透过某个人类群体殷勤奉送的、有专利的眼镜去看——那群体或大或小,出于这样那样的理由被赋予了权威。

这种“倾向”,或这种内心的态度,不是成熟的特权,而是童年的特权。它与生命一同,是出生时得到的馈赠——一份谦卑而可畏的馈赠。一份常常被深埋的馈赠,有些人还能多少保住它,或者也许能重新找回它……

也可以把它叫作孤独的馈赠。
